\chapter{The Forward Transliterator}

\section{From words to segments}
For each word the engine obtains a phoneme list, aligns vowel phonemes to the vowel letter groups of the spelling when their counts agree, and builds a list of segments. Each segment is a letter with a kind (consonant, long vowel letter, glide, seat) and an optional mark.

Consonant phonemes append their letters, digraphs such as \ar{تْش} appending two segments of one unit. Vowel phonemes attach their mark to the last consonant, or to a preceding \ar{ي}, or to a new hamza seat, and then append their vowel letters. Rhotic vowels append \ar{ر}.

\section{The sukoon pass}
After all segments are placed, every consonant or glide without a vowel receives sukoon. Long vowel letters and seats do not. The resulting form is the full-density spelling.

\section{Mark density}
A final filter produces the other densities, so density can never change a letter.
\begin{description}
\item[full] all marks, including sukoon on every vowelless consonant.
\item[medial] sukoon removed from the end of the word.
\item[vowels] fatha, damma and kasra only.
\item[none] no marks. This is the consonantal skeleton with vowel letters, the form requested in the earliest experiments.
\end{description}

\section{Attested forms}
With \code{--attested} the engine prefers a form recorded in a specimen whenever the word occurs there. The list is built by aligning each specimen token to its English word and asserting that token counts match (\code{scripts/build\_attested.py}). Page forms are listed before typed forms.

\section{Demonstration}
Table~\ref{tab:demo} shows one sentence of the page under several settings. The first row is the page itself.

\begin{table}[ht]\centering\small
\begin{tabular}{@{}lp{0.72\textwidth}@{}}
\toprule
Setting & Output\\
\midrule
full marks & \ar{إِنْ ذِسْ بَارْتْ أُفْ هِيلِنْ رُولِنْزْ سَايْكُوسِينِمَا، شِي دِلْفْزْ إِنْتُو ذَا نَيْتْشَر أُفْ لَاكْ.}\\[3pt]
computed, full & \ar{إِنْ ذِسْ بَارْتْ أُفْ هِيلِنْ رُولِنْزْ سَايْكُوسِينِمَا، شِي دِلْفْزْ إِنْتُو ذَا نَيْتْشَرْ أُفْ لَاكْ.}\\[3pt]
medial sukoon & \ar{إِن ذِس بَارْت أُف هِيلِن رُولِنْز سَايْكُوسِينِمَا، شِي دِلْفْز إِنْتُو ذَا نَيْتْشَر أُف لَاك.}\\[3pt]
vowel marks only & \ar{إِن ذِس بَارت أُف هِيلِن رُولِنز سَايكُوسِينِمَا، شِي دِلفز إِنتُو ذَا نَيتشَر أُف لَاك.}\\[3pt]
unmarked & \ar{إن ذس بارت أف هيلن رولنز سايكوسينما، شي دلفز إنتو ذا نيتشر أف لاك.}\\[3pt]
v as feh, g as jeem & \ar{إِنْ ذِسْ بَارْتْ أُفْ هِيلِنْ رُولِنْزْ سَايْكُوسِينِمَا، شِي دِلْفْزْ إِنْتُو ذَا نَيْتْشَرْ أُفْ لَاكْ.}\\[3pt]
\bottomrule
\end{tabular}
\caption{One sentence in several settings. The second row is the computed form with the defaults.}\label{tab:demo}
\end{table}

\section{The density ladder}
The four densities are best seen on a single sentence. Each line contains the same letters; only the marks differ.

\begin{table}[ht]\centering\small
\begin{tabular}{@{}ll@{}}
\toprule
Density & Result\\
\midrule
full & \ar{ذَا لِتَرْزْ تْشَيْنْجْ، بَتْ ذَا وُرْدْزْ رِمَيْنْ.}\\[3pt]
medial & \ar{ذَا لِتَرْز تْشَيْنْج، بَت ذَا وُرْدْز رِمَيْن.}\\[3pt]
vowels & \ar{ذَا لِتَرز تشَينج، بَت ذَا وُردز رِمَين.}\\[3pt]
none & \ar{ذا لترز تشينج، بت ذا وردز رمين.}\\[3pt]
\bottomrule
\end{tabular}
\caption{One sentence at each density.}\label{tab:density}
\end{table}

\section{Worked passages}
The following passages are public-domain texts and short originals, set with the default options. Each block shows the Arabic-script form, a Latin respelling that is only a reading aid, and the English. They illustrate how the rules behave on running prose and not on isolated words, and they include some of the weaknesses that Chapter~\ref{ch:open} lists.

\subsection{Children's prose}
\begin{quote}\noindent \ar{أَلِسْ وَازْ بِغِنِنْغْ تُو غِتْ فِيرِي تَايَرْدْ أُفْ سِتِنْغْ بَايْ هِرْ سِسْتَرْ أُونْ ذَا بَانْغْكْ، أَنْدْ أُفْ هَافِنْغْ نُوثِنْغْ تُو دُو.}\\[2pt]
\ar{ألس واز بغننغ تو غت فيري تايرد أف ستنغ باي هر سستر أون ذا بانغك، أند أف هافنغ نوثنغ تو دو.}\\[1pt]
\textit{alus waaz bigining too get veree tayerd uv siting bay her sister aan dhu bangk und uv having nuthing too doo}\\[1pt]
Alice was beginning to get very tired of sitting by her sister on the bank, and of having nothing to do.
\end{quote}
\begin{quote}\noindent \ar{وُونْسْ أُورْ تْوَايْسْ شِي هَادْ بِيبْتْ إِنْتُو ذَا بُوكْ هِرْ سِسْتَرْ وَازْ رِيدِنْغْ، بَتْ إِتْ هَادْ نُو بِكْتْشَرْزْ أُورْ كُونْفَرْسَيْشَنْزْ إِنْ إِتْ.}\\[2pt]
\ar{وونس أور توايس شي هاد بيبت إنتو ذا بوك هر سستر واز ريدنغ، بت إت هاد نو بكتشرز أور كونفرسيشنز إن إت.}\\[1pt]
\textit{wuns awr tways shee had peept intoo dhu buk her sister waaz reeding but it had noh pikcherz awr kaanverseyshunz in it}\\[1pt]
Once or twice she had peeped into the book her sister was reading, but it had no pictures or conversations in it.
\end{quote}
\begin{quote}\noindent \ar{أَنْدْ وَتْ إِزْ ذَا يُوسْ أُفْ أَ بُوكْ، ثُوتْ أَلِسْ، وِثَاوْتْ بِكْتْشَرْزْ أُورْ كُونْفَرْسَيْشَنْزْ؟}\\[2pt]
\ar{أند وت إز ذا يوس أف أ بوك، ثوت ألس، وثاوت بكتشرز أور كونفرسيشنز؟}\\[1pt]
\textit{und wut iz dhu yoos uv u buk thawt alus withowt pikcherz awr kaanverseyshunz}\\[1pt]
And what is the use of a book, thought Alice, without pictures or conversations?
\end{quote}

\subsection{Scripture}
\begin{quote}\noindent \ar{إِنْ ذَا بِغِنِنْغْ غُودْ كْرَيْتَدْ ذَا هِيفِنْ أَنْدْ ذَا أَرْثْ.}\\[1pt]
\textit{in dhu bigining gaad kreeeytud dhu hevun und dhu erth}\\[1pt]
In the beginning God created the heaven and the earth.
\end{quote}
\begin{quote}\noindent \ar{أَنْدْ ذَا أَرْثْ وَازْ وِثَاوْتْ فُورْمْ، أَنْدْ فُويْدْ؛ أَنْدْ دَارْكْنِسْ وَازْ أَبُونْ ذَا فَيْسْ أُفْ ذَا دِيبْ.}\\[1pt]
\textit{und dhu erth waaz withowt fawrm und voyd und daarknus waaz upaan dhu feys uv dhu deep}\\[1pt]
And the earth was without form, and void; and darkness was upon the face of the deep.
\end{quote}
\begin{quote}\noindent \ar{أَنْدْ غُودْ سِدْ، لِتْ ذِرْ بِي لَايْتْ: أَنْدْ ذِرْ وَازْ لَايْتْ.}\\[1pt]
\textit{und gaad sed let dher bee layt und dher waaz layt}\\[1pt]
And God said, Let there be light: and there was light.
\end{quote}

\subsection{Verse}
\begin{quote}\noindent \ar{شَالْ آيْ كُمْبِرْ ذِي تُو أَ سَمَرْزْ دَيْ؟}\\[1pt]
\textit{shal ay kumper dhee too u sumerz dey}\\[1pt]
Shall I compare thee to a summer's day?
\end{quote}
\begin{quote}\noindent \ar{ذَاوْ آرْتْ مُورْ لَفْلِي أَنْدْ مُورْ تِمْبْرَتْ.}\\[1pt]
\textit{dhow aart mawr luvlee und mawr temprut}\\[1pt]
Thou art more lovely and more temperate.
\end{quote}
\begin{quote}\noindent \ar{رُوفْ وِنْدْزْ دُو شَيْكْ ذَا دَارْلِنْغْ بَدْزْ أُفْ مَيْ، أَنْدْ سَمَرْزْ لِيسْ هَاثْ أُولْ تُو شُورْتْ أَ دَيْتْ.}\\[1pt]
\textit{ruf windz doo sheyk dhu daarling budz uv mey und sumerz lees hath awl too shawrt u deyt}\\[1pt]
Rough winds do shake the darling buds of May, and summer's lease hath all too short a date.
\end{quote}
\begin{quote}\noindent \ar{وُونْسْ أَبُونْ أَ مِدْنَايْتْ دْرِيرِي، وَايْلْ آيْ بَانْدَرْدْ، وِيكْ أَنْدْ وِيرِي،}\\[1pt]
\textit{wuns upaan u midnayt driree wayl ay paanderd week und wiree}\\[1pt]
Once upon a midnight dreary, while I pondered, weak and weary,
\end{quote}
\begin{quote}\noindent \ar{أُوفَرْ مِنِي أَ كْوَيْنْتْ أَنْدْ كْيُرِيَسْ فُولْيُومْ أُفْ فَرْغُوتِنْ لُورْ.}\\[1pt]
\textit{ohver menee u kweynt und kyureeus vaalyoom uv fergaatun lawr}\\[1pt]
Over many a quaint and curious volume of forgotten lore.
\end{quote}

\subsection{A long sentence}
\begin{quote}\noindent \ar{وِنْ إِنْ ذَا كُورْسْ أُفْ هْيُومَنْ إِفِنْتْسْ، إِتْ بِكَمْزْ نِيسِسِرِي فُورْ وُونْ بِيبَلْ تُو دِزُولْفْ ذَا بُلِيتِكَلْ بَانْدْزْ وِتْشْ هَافْ كُنِكْتِدْ ذِمْ وِذْ أَنُوذَرْ، أَنْدْ تُو أَسُومْ أَمُونْغْ ذَا بَاوَرْزْ أُفْ ذَا أَرْثْ، ذَا سِيبَرَيْتْ أَنْدْ إِيكْوَلْ سْتَيْشَنْ تُو وِتْشْ ذَا لُوزْ أُفْ نَيْتْشَرْ أَنْدْ أُفْ نَيْتْشَرْزْ غُودْ إِنْتَايْتَلْ ذِمْ، أَ دِيسِنْتْ رِسْبِكْتْ تُو ذَا أُبِينْيَنْزْ أُفْ مَانْكَايْنْدْ رِكْوَايَرْزْ ذَتْ ذَيْ شُودْ دِكْلِرْ ذَا كَازِزْ وِتْشْ إِمْبِلْ ذِمْ تُو ذَا سِيبَرَيْشَنْ.}\\[1pt]
\textit{wen in dhu kawrs uv hyoomun ivents it bikumz nesuseree fawr wun peepul too dizaalv dhu pulitukul bandz wich hav kunektid dhem widh unudher und too usoom umung dhu powerz uv dhu erth dhu sepereyt und eekwul steyshun too wich dhu lawz uv neycher und uv neycherz gaad entaytul dhem u deesunt rispekt too dhu upinyunz uv mankaynd reekwayerz dhat dhey shud dikler dhu kaazuz wich impel dhem too dhu sepereyshun}\\[1pt]
When in the Course of human events, it becomes necessary for one people to dissolve the political bands which have connected them with another, and to assume among the powers of the earth, the separate and equal station to which the Laws of Nature and of Nature's God entitle them, a decent respect to the opinions of mankind requires that they should declare the causes which impel them to the separation.
\end{quote}

\subsection{Narrative}
\begin{quote}\noindent \ar{ذَا تْرَافِلَرْ كَيْمْ دَاوْنْ فْرُومْ ذَا سْتَارْ سِرَسْ تُو آوَرْ لِتَلْ أَنْتْ هِلْ، أَنْدْ مِزَرْدْ ذَا أَرْثْ وِذْ هِزْ فِنْغَرْزْ.}\\[1pt]
\textit{dhu travuler keym down frum dhu staar sirius too ower litul ant hil und mezherd dhu erth widh hiz finggerz}\\[1pt]
The traveller came down from the star Sirius to our little ant hill, and measured the earth with his fingers.
\end{quote}
\begin{quote}\noindent \ar{هِي فَاوْنْدْ ذَا بِيبَلْ سْمُولْ، بَتْ هِي فَاوْنْدْ ذِمْ كْيُرِيَسْ، أَنْدْ هِي أَسْكْتْ ذِمْ مِنِي كْوِسْتْشَنْزْ أَبَاوْتْ ذِرْ لِفْزْ.}\\[1pt]
\textit{hee fownd dhu peepul smawl but hee fownd dhem kyureeus und hee askt dhem menee kweschunz ubowt dher livz}\\[1pt]
He found the people small, but he found them curious, and he asked them many questions about their lives.
\end{quote}

\section{Names, dates and numbers}
Proper names are looked up like other words, and the fallback handles names the dictionary lacks. Numbers can be written in Arabic-Indic digits, in Latin digits or spoken, depending on the \texttt{digits} option. The spoken form is useful for text meant to be read aloud.

\begin{quote}\noindent \ar{هِيلِنْ رُولِنْزْ لِفْزْ إِنْ هَالِفَاكْسْ، أَنْدْ هِرْ بْرُوذَرْ فِكْتَرْ لِفْزْ إِنْ فَنْكُوفَرْ.}\\[1pt]
\textit{helun raalinz livz in halifaks und her brudher vikter livz in vankoover}\\[1pt]
Helen Rollins lives in Halifax, and her brother Victor lives in Vancouver.
\end{quote}
\begin{quote}\noindent \ar{تُوزْدِي، ذَا فِفْثْ أُفْ أُوكْتُوبَرْ، أَتْ سِيفِنْ ثِرْدِي، ذَا تْرَيْنْ فْرُومْ تَرُونْتُو أَرَايْفْزْ.}\\[1pt]
\textit{toozdee dhu fifth uv aaktohber at sevun therdee dhu treyn frum teraantoh erayvz}\\[1pt]
Tuesday, the fifth of October, at seven thirty, the train from Toronto arrives.
\end{quote}
\begin{quote}\noindent \ar{ذِرْ أَرْ ٢٥ وُرْدْزْ أُونْ ذِسْ بَيْجْ أَنْدْ ٣ أُفْ ذِمْ أَرْ نَمْبَرْزْ.}\\[1pt]
\textit{dher aar werdz aan dhis peyj und uv dhem aar numberz}\\[1pt]
There are 25 words on this page and 3 of them are numbers.
\end{quote}
\begin{quote}\noindent \ar{إِنْ ٢٠٢٦ إِتْ كُوسْتْ ١,٢٥٠ دُولَرْزْ.}\\[1pt]
\textit{in it kaast daalerz}\\[1pt]
In 2026 it cost 1,250 dollars.
\end{quote}

\section{Reading the results}
Three habits of the output are worth noting. Common words always take their fixed form, so \emph{the}, \emph{of}, \emph{and} and \emph{is} look the same everywhere. Stressed vowels are usually marked more heavily than unstressed ones, which gives longer words a visible rhythm. And function words with weak forms in speech, such as \emph{to} and \emph{a}, keep their fixed spellings and are not reduced.
